\chapter{Eye Floaters (Vitreous Floaters)}

\section*{Definition}
Eye floaters are small specks, spots, or thread-like shapes that drift across the field of vision. They are most noticeable when looking at a plain bright background. While usually benign, the sudden onset of floaters, especially with flashes, requires urgent evaluation.

\section*{What Happens in Your Body}
Floaters are shadows cast on the retina by clumps within the vitreous humor of the eye. With age, the vitreous humor undergoes syneresis (liquefaction) and contraction, causing collagen fibers to aggregate. Posterior vitreous detachment (PVD) can pull on the retina, causing the sudden appearance of floaters and flashes.

\section*{Causes}
\begin{itemize}
\item Age-related vitreous degeneration (most common)
\item Posterior vitreous detachment
\item Retinal tear
\item Retinal detachment
\item Vitreous hemorrhage (diabetic retinopathy, trauma)
\item Inflammatory cells (uveitis)
\item Myopia (nearsightedness)
\item Eye trauma or surgery
\item Migraine (visual aura with scintillating scotomas)
\end{itemize}

\section*{When to See a Doctor}
\begin{itemize}
\item Small spots, specks, or cobweb-like shapes in vision
\item Floaters that drift slowly across the visual field
\item Shadows or dark shapes that move with eye movement
\item Sudden increase in number or size of floaters
\item Flashes of light in peripheral vision
\item Curtain-like shadow or loss of vision
\end{itemize}

\section*{Related Symptoms}
\begin{itemize}
\item Flashing lights (photopsia)
\item Vision loss or visual field defect
\item Blurred vision
\item Retinal detachment
\item Posterior vitreous detachment
\item Migraine with aura
\end{itemize}

\section*{Self-Care / Home Management}
\begin{itemize}
\item Move the eyes up and down to shift the floater out of the visual axis
\item Try to ignore benign floaters (they often become less noticeable over time)
\item Seek immediate evaluation if new floaters or flashes appear suddenly
\item Do not rub the eyes
\item Keep a healthy diet rich in antioxidants
\end{itemize}

\section*{How Your Doctor Will Evaluate You}
\begin{itemize}
\item Comprehensive eye examination with pupillary dilation
\item Slit-lamp biomicroscopy of the vitreous
\item Indirect ophthalmoscopy of the retina
\item B-scan ultrasound if view is obstructed
\item Optical coherence tomography (OCT) of the retina
\end{itemize}

\section*{Treatment Options}
\begin{itemize}
\item Observation and reassurance for benign floaters
\item Treatment of underlying cause (diabetic retinopathy, uveitis)
\item Laser vitreolysis (YAG laser for symptomatic floaters)
\item Vitrectomy for severe, disabling floaters (rarely needed)
\item Surgical repair for retinal tear or detachment
\end{itemize}
\section*{References}
\begin{thebibliography}{99}
\bibitem{ref1} Hollands H, Johnson D, Brox AC, et al. "Acute-Onset Floaters and Flashes: Is This Patient at Risk for Retinal Detachment?" JAMA. 2009;302(20):2243--2249.
\bibitem{ref2} Sebag J. "Floaters and the Quality of Life." Am J Ophthalmol. 2011;152(1):3--4.
\bibitem{ref3} Milston R, Madigan MC, Sebag J. "Vitreous Floaters: Etiology, Diagnostics, and Management." Surv Ophthalmol. 2016;61(2):211--227.
\bibitem{ref4} Lumi X, Hawlina M, Petrovski G, et al. "Vitreous Floaters: A Review of Current Management." Int J Ophthalmol. 2020;13(7):1135--1141.
\bibitem{ref5} Johnson MW. "Posterior Vitreous Detachment." In: Ryan SJ, ed. Retina. 6th ed. Elsevier; 2018.
\end{thebibliography}
